\chapter{Alte und neue Theologie}
\quellseite{31}

Ein neuer Theologe meint von seinem Rathgeber aus in das Innere der hl. Dreieinigkeit geschaut und die persönlichen Eigenschaften derselben begriffen zu haben, und sagt, daß er
„\gesperrt{zu figürlicher Versinnbildlichung des trinitarischen Processes würde einen feurigen Kreis malen als Sinnbild des feurigen Liebeswillens des Vaters, und in diesem Kreise ein sonnenlichtes Centrum als Sinnbild des Sohnes, des die ganze unendliche Tiefe des göttlichen Liebeswesens lichtenden Liebeszieles, und von diesem sonnenlichten Centrum ausgehend bis zur Peripherie des feurigen Kreises eine Strahlenfülle als Sinnbild des vom Vater durch den Sohn hindurch ausgehenden und}
\quellseite{32}
\gesperrt{Vater und Sohn erfüllenden Liebestriumphes}.“ Vgl. „L. u. W.“, 2. Bd., Januarheft pag. 24. — Ein alter Theologe (Gregor von Nazianz, gest. 390) dagegen sagt: „Von Gottes Geburt soll man mit Ehrfurchtigkeit still schweigen; halte es bei dir für ein großes, daß du so viel gelernt, daß Gottes Sohn geboren sei. Wie er aber geboren sei, das ist auch über aller Engel Verstand, geschweige daß wir Menschen es verstehen sollten. Was ist aber der Ausgang des Hl. Geistes? Sag’ mir erst, wie es mit des Vaters Eigenschaft bewandt sei, nach welcher Er nicht geboren worden, so will ich dir alsdann sagen, wie der Sohn geboren worden sei, und wie der Hl. Geist ausgehe. Dergestalt werden wir alle beide zu Thoren werden, wenn wir in die Geheimnisse Gottes hineinsehen wollen.“ B.
